\section{La membrana plasmatica: funzioni}

% ---------------------------------------------------------------------------
\subsection{Funzioni generali}

\begin{itemize}
  \item protettiva;
  \item forma;
  \item scambio di sostanze nutritive;
  \item trasporto;
  \item reattività agli stimoli esterni;
  \item risposta agli ormoni;
  \item interazione e comunicazione tra le cellule;
  \item proprietà antigeniche e reazioni immunitarie.
\end{itemize}

% ---------------------------------------------------------------------------
\subsection{Differenziazione funzionale (cellula epiteliale)}

\begin{itemize}
  \item \textbf{membrana apicale}: regolazione dell'ingresso di materiali nutritizi e di acqua; regolazione della secrezione; protezione;
  \item \textbf{membrana laterale}: contatto e adesione fra cellule; comunicazione cellulare;
  \item \textbf{membrana basale}: contatto con il substrato; creazione di gradienti ionici.
\end{itemize}

% ---------------------------------------------------------------------------
\subsection{Funzioni delle proteine di membrana}

\begin{itemize}
  \item \textbf{ancoraggio}: alcune proteine (es.\ le integrine) ancorano la cellula alla matrice extracellulare e si connettono ai microfilamenti intracellulari;
  \item \textbf{trasporto passivo}: certe proteine formano canali per il passaggio selettivo di ioni o molecole;
  \item \textbf{trasporto attivo}: alcune proteine pompano i soluti attraverso la membrana, con apporto diretto di energia;
  \item \textbf{attività enzimatica}: molti enzimi legati alla membrana catalizzano reazioni all'interno o sulla superficie;
  \item \textbf{trasduzione del segnale}: alcuni recettori legano molecole segnale (es.\ ormoni) e trasmettono l'informazione all'interno della cellula;
  \item \textbf{riconoscimento cellulare}: alcune glicoproteine fungono da marcatori di identificazione (es.\ gli antigeni di superficie delle cellule batteriche, riconosciuti come estranei);
  \item \textbf{giunzione intercellulare}: le proteine di adesione legano le membrane di cellule adiacenti.
\end{itemize}

% ---------------------------------------------------------------------------
\subsection{Diffusione e osmosi}

\rubrica{Diffusione}
Movimento delle molecole da una zona a maggiore concentrazione a una a minore concentrazione, fino a distribuzione uniforme (es.\ una zolletta di zucchero immersa in acqua si dissolve e diffonde uniformemente).

\rubrica{Osmosi}
Movimento netto di acqua attraverso una membrana selettivamente permeabile (permeabile all'acqua ma non al soluto). La \textbf{pressione osmotica} della soluzione è pari alla forza che occorre applicare per impedire l'ascesa del livello del fluido.
\begin{itemize}
  \item soluzione \textbf{ipotonica}: guadagno netto di acqua $\rightarrow$ la cellula si gonfia;
  \item soluzione \textbf{ipertonica}: perdita di acqua $\rightarrow$ la cellula si raggrinzisce;
  \item soluzione \textbf{isotonica}: né guadagno né perdita di acqua.
\end{itemize}

% ---------------------------------------------------------------------------
\subsection{Meccanismi di trasporto attraverso la membrana}

\begin{itemize}
  \item \textbf{trasporto passivo} (secondo gradiente, senza energia):
    \begin{itemize}
      \item non mediato: diffusione semplice (anche attraverso un canale acquoso);
      \item mediato: diffusione facilitata secondo gradiente di concentrazione;
    \end{itemize}
  \item \textbf{trasporto attivo}: contro gradiente, con apporto diretto di energia.
\end{itemize}

\rubrica{Trasporto passivo: diffusione semplice}
Piccole molecole (es.\ $O_2$, $CO_2$) attraversano direttamente il doppio strato lipidico.

\rubrica{Trasporto passivo: diffusione facilitata}
\begin{itemize}
  \item \textbf{proteine canale}: formano un poro selettivo (es.\ acquaporina, canale per l'acqua);
  \item \textbf{proteine carrier} (trasportatori): legano il soluto e cambiano conformazione esponendo il sito di legame alternativamente verso l'esterno o verso l'interno (legame $\rightarrow$ trasporto $\rightarrow$ dissociazione $\rightarrow$ ritorno alla conformazione iniziale); es.\ il trasportatore del glucosio GLUT1.
\end{itemize}

\rubrica{Trasporto attivo}
Contro gradiente, con consumo di energia (ATP).
\begin{itemize}
  \item \textbf{pompa $Na^+$--$K^+$}: ciclo conformazionale che, idrolizzando ATP, espelle $Na^+$ e importa $K^+$;
  \item \textbf{trasporto attivo secondario}: $Na^+$ e glucosio si legano insieme alla proteina carrier; la proteina cambia conformazione e li rilascia nel citosol. Il glucosio è trasportato contro il proprio gradiente sfruttando il gradiente del sodio.
\end{itemize}

% ---------------------------------------------------------------------------
\subsection{Trasporto di massa: endocitosi ed esocitosi}

\rubrica{Esocitosi}
Una vescicola secretoria si avvicina alla membrana plasmatica e si fonde con essa $\rightarrow$ rilascia il proprio contenuto nell'ambiente extracellulare; le proteine inserite nella membrana della vescicola diventano parte integrante della membrana plasmatica.

\rubrica{Fagocitosi}
Pieghe (pseudopodi) della membrana circondano la particella da ingerire formando un vacuolo (fagosoma) $\rightarrow$ il vacuolo, in seguito a una strozzatura, si libera all'interno della cellula $\rightarrow$ i lisosomi si fondono con esso (fagolisosoma) e riversano i loro enzimi digestivi sul materiale ingerito. Passaggi: intrappolamento $\rightarrow$ inglobamento $\rightarrow$ digestione $\rightarrow$ assorbimento. Nei mammiferi alcuni globuli bianchi (\textit{fagociti}) svolgono un processo simile.

\rubrica{Pinocitosi (endocitosi di massa)}
Goccioline di fluido sono intrappolate da pieghe della membrana plasmatica $\rightarrow$ subiscono una strozzatura e divengono piccole vescicole piene di fluido $\rightarrow$ il contenuto viene trasferito al citosol.

\rubrica{Endocitosi mediata da recettore}
Specifiche macromolecole si legano a recettori situati nelle fossette rivestite (di clatrina). Esempio: assorbimento delle LDL.
\begin{enumerate}
  \item la LDL si attacca a recettori specifici nelle fossette rivestite della membrana plasmatica;
  \item l'endocitosi forma una vescicola rivestita nel citosol; il rivestimento viene subito rimosso;
  \item la vescicola non rivestita si fonde con un endosoma;
  \item i recettori vengono riciclati e tornano sulla membrana plasmatica;
  \item la vescicola contenente le LDL si fonde con un lisosoma (lisosoma secondario); gli enzimi idrolitici rilasciano dalle LDL il colesterolo, utilizzato dalla cellula.
\end{enumerate}
